\section{The quintic specialization: a completely explicit Dirichlet descent}
\label{sec:five}

This section proves the exponent-five specialization, including both parity
cases.  Its arithmetic takes place in the fixed quadratic ring
\[
 R=\mathbb Z[\varphi],\qquad
 \varphi^2=\varphi+1,\qquad
 \varphi=\frac{1+\sqrt5}{2}.
\]
The square-root notation is an abbreviation for integer-pair calculations,
not an assumption that iPRA contains real numbers.  The proof below replaces
the usual appeal to a least counterexample by a bounded iteration.

\subsection{Arithmetic in the golden quadratic ring}

Represent \(a+b\varphi\) by the integer pair \((a,b)\).  Addition is
coordinatewise, and multiplication, conjugation, and norm are
\[
 (a,b)(c,d)=(ac+bd,\,ad+bc+bd),\qquad
 \overline{(a,b)}=(a+b,-b),\qquad
 N(a,b)=a^2+ab-b^2.
\]
Direct expansion proves
\(\alpha\bar\alpha=N(\alpha)\) and
\(N(\alpha\beta)=N(\alpha)N(\beta)\).
The norm of a nonzero pair is nonzero: otherwise
\((2a+b)^2=5b^2\), and cancelling the gcd of \(2a+b\) and \(b\)
contradicts the fact that a prime dividing a square divides its base.
All this uses only integer arithmetic.

For a convenient presentation of inequalities we use the ordered quadratic
field \(\mathbb Q(\sqrt5)\).  It too is an arithmetic coding device:
the sign of \(A+B\sqrt5\), after clearing positive denominators, is decided
by the signs of \(A,B\), and, when these have opposite signs, by comparing
\(A^2\) with \(5B^2\).  Equality in the latter comparison is impossible
unless \(A=B=0\).  The ordered-field laws follow by multiplying out these
finite sign cases.  In particular, no principle about arbitrary real
numbers, completeness, or choice is used.  The elementary inequalities
\[
 2<\sqrt5<3,\qquad \frac32<\varphi<2,\qquad
 \varphi^{-1}=\varphi-1
\]
are verified by squaring positive rational numbers.

\paragraph{Euclidean division.}
Write \(\alpha/\beta=x+y\varphi\) with rational \(x,y\) and
\(\beta\ne0\).  Choose nearest integers \(m,n\) by quotient and remainder,
breaking ties in a fixed direction.  Put
\(r=x-m\), \(s=y-n\).  Since \(|r|,|s|\le1/2\),
\[
 |N(r+s\varphi)|
 =|r^2+rs-s^2|\le r^2+|rs|+s^2\le\frac34.
\]
Consequently
\[
 \rho=\alpha-(m+n\varphi)\beta
 \quad\Longrightarrow\quad
 |N(\rho)|\le\frac34|N(\beta)|<|N(\beta)|.
\]
Thus the ordinary Euclidean algorithm works with the nonnegative integer
measure \(|N|\).  If its initial second argument has norm magnitude \(M\),
there are at most \(M+1\) divisions.  Back substitution gives a Bezout
identity for its gcd.

Here are the elementary consequences, with their constructive content
made explicit.  An element is a unit exactly when its norm is \(1\) or
\(-1\), since then \(\bar\alpha/N(\alpha)\) is its inverse.  An irreducible
element \(\pi\) is prime: if \(\pi\nmid a\), its gcd with \(a\) is a unit,
so a Bezout identity \(u\pi+va=1\), multiplied by \(b\), proves
\(\pi\mid b\) whenever \(\pi\mid ab\).  Factorization terminates because
a factorization into two nonunits splits the positive integer \(|N|\)
into two factors each at least \(2\).  Uniqueness follows by repeatedly
using this prime-divisor property and cancelling.

It is useful to explain why this argument need not conceal unbounded
search.  Let \(\alpha=a+b\varphi\ne0\) and
\(H=|a|+|b|+1\).  In the positive embedding, both
\(|\alpha|\) and \(1/|\alpha|\) are less than \(2H\): the first
bound is immediate, and the second follows from
\(|N(\alpha)|\ge1\) and \(|\bar\alpha|<2H\).
The inequality
\(\varphi^j\ge1+j/2\), proved by ordinary induction, shows that a
signed exponent \(k\) with \(|k|\le4H\) normalizes a signed associate
\[
 1\le v=\pm\alpha\varphi^k<\varphi.
\]
Explicitly, the endpoints \(\varphi^{-4H}|\alpha|<1\) and
\(\varphi^{4H}|\alpha|>1\) bracket the crossing of \(1\).
Scanning consecutive exponents and taking the first value at least \(1\)
gives the required strict upper bound \(\varphi\).
If \(M=|N(\alpha)|\), then \(|\bar v|\le M\).  Solving for the two
integer coordinates of \(v\) bounds both by \(3M+6\).
Therefore, if \(\alpha\) is reducible, a nonunit proper divisor can be
found among the finitely many pairs in the box
\([-(3M+6),3M+6]^2\), by checking divisibility and
\(2\le |N(\delta)|<M\).
Each quotient is computed by the conjugate formula and a test that its
two rational coordinates are integers.  This supplies a bounded
factorization procedure; at most \(M\) nonunit factors are needed.
One can normalize factors as they are produced.  Thus the preceding
existence and uniqueness arguments have primitive-recursive
implementations, with quantifier-free invariants on finite lists.

\paragraph{All units.}
Every unit is \(\pm\varphi^j\), \(j\in\mathbb Z\).
For completeness, normalize a unit as above to \(1\le v<\varphi\),
and write \(v=a+b\varphi\).  If \(N(v)=1\), then
\[
 0\le b=\frac{v-v^{-1}}{\sqrt5}
 <\frac{\varphi-\varphi^{-1}}{\sqrt5}
 =\frac1{\sqrt5}<1.
\]
Thus \(b=0\), and \(v=1\).  If \(N(v)=-1\), then
\[
 0<b=\frac{v+v^{-1}}{\sqrt5}<1,
\]
which is impossible.  For the last strict inequality use
\(\varphi+\varphi^{-1}=\sqrt5\) and the identity
\[
 (\varphi+\varphi^{-1})-(v+v^{-1})
   =(\varphi-v)\left(1-\frac1{\varphi v}\right)>0.
\]
Undoing the normalization proves the assertion.
In particular, every unit is \(u^5\varphi^r\) for a unit \(u\)
and a uniquely selected \(r\in\{0,1,2,3,4\}\);
the sign is absorbed in a fifth power and the exponent is divided by \(5\).

\subsection{A fifth-power extraction lemma}

We need to remove the unit in a coprime factorization.  Reduction modulo
\(5\) is especially simple:
\[
 \varphi^5=3+5\varphi,\qquad
 (a+b\varphi)^5\equiv a+3b\pmod{5R}.
\]
Thus a fifth power is a rational scalar modulo \(5R\).  The coordinate
pairs of \(1,\varphi,\ldots,\varphi^4\) are
\[
 (1,0),\quad(0,1),\quad(1,1),\quad(1,2),\quad(2,3).
\]
Suppose \(\alpha=A+B\varphi\), \(5\mid B\), \(5\nmid A\), and
\(\alpha\bar\alpha=c^5\), where \(\alpha,\bar\alpha\) are coprime in
\(R\).  Unique factorization shows that each prime multiplicity in
\(\alpha\) is a multiple of \(5\), since the prime occurs in no factor
of \(\bar\alpha\).  Hence
\(\alpha=\varepsilon\delta^5=\varphi^r\gamma^5\).
Modulo \(5R\), \(\gamma^5\) is a nonzero scalar: a zero scalar
would make both coordinates of \(\alpha\) divisible by \(5\).
The displayed table and \(5\mid B\) now force \(r=0\).
We have proved
\[
 \boxed{\quad
 \alpha\bar\alpha=c^5,\quad
 (\alpha,\bar\alpha)=1,\quad
 5\mid B,\quad5\nmid A
 \quad\Longrightarrow\quad \alpha=\gamma^5.
 \quad}
 \tag{Q1}\label{eq:quintic-extraction}
\]

Set
\[
 F(t,s)=t^4+10t^2s^2+5s^4,\qquad
 G(t,s)=t^4+50t^2s^2+125s^4.
\]
Expansion gives
\[
 (t+s\sqrt5)^5=tG(t,s)+5sF(t,s)\sqrt5.
 \tag{Q2}\label{eq:quintic-expansion}
\]
We record the two precise forms of (Q1) that will be used.

\begin{lemma}[Direct coordinates]\label{lem:quintic-direct}
Let \(P,Q,c\) be positive integers with
\[
 (P,Q)=1,\quad P\text{ odd},\quad Q\text{ even},\quad
 5\mid Q,\quad5\nmid P,\quad P^2-5Q^2=c^5.
\]
Then there are positive coprime integers \(t,s\), with \(t\) odd,
\(s\) even and \(5\nmid t\), such that
\[
 P=tG(t,s),\qquad Q=5sF(t,s).
 \tag{Q3}\label{eq:quintic-direct}
\]
\end{lemma}

\begin{proof}
Put \(\alpha=P+Q\sqrt5=(P-Q)+2Q\varphi\).
The integer \(\gcd(P^2-5Q^2,2P)\) is \(1\):
an odd prime common divisor would divide \(P\) and \(5Q^2\),
contrary to \((P,Q)=1\) and \(5\nmid P\), and the norm is odd.
A Bezout identity between the product
\(\alpha\bar\alpha=P^2-5Q^2\) and the sum
\(\alpha+\bar\alpha=2P\) gives a Bezout identity between
\(\alpha\) and \(\bar\alpha\).  Indeed
\(u\alpha\bar\alpha+v(\alpha+\bar\alpha)=1\) becomes
\((u\bar\alpha+v)\alpha+v\bar\alpha=1\).
The coordinate hypotheses of (Q1) hold, so \(\alpha=\gamma^5\).

Write \(\gamma=a+b\varphi\).  Modulo \(2R\) one has
\(\gamma^4=\gamma\), by checking the four coordinate pairs, and hence
\(\gamma^5=\gamma^2\).  The second coordinate of
\(\gamma^2\) is \(b^2\equiv b\pmod2\).
The second coordinate of \(\alpha\) is \(2Q\), so \(b\) is even.
Thus \(\gamma=t+s\sqrt5\) with integers
\(s=b/2\), \(t=a+b/2\).
Its norm is odd, so \(t,s\) have opposite parity.

Formula (Q2) proves (Q3).  Its two coefficient polynomials have positive
even-power terms.  Since \(P,Q>0\), both \(t,s\) are positive.
If \(t\) were even and \(s\) odd, \(5sF(t,s)\) would be odd,
contrary to the parity of \(Q\).  Thus \(t\) is odd and \(s\) even.
The first equation modulo \(5\) gives \(5\nmid t\).

Finally, \(\gamma,\bar\gamma\) are coprime: a Bezout identity between
their fifth powers is also a Bezout identity between them after
absorbing fourth powers into the coefficients.  If a rational prime
divided both \(t,s\), it would divide both \(\gamma,\bar\gamma\);
that prime is a nonunit of \(R\), since its norm is its square.
This is impossible.  Hence \((t,s)=1\).
\end{proof}

\begin{lemma}[Half-integral coordinates]\label{lem:quintic-half}
Let \(P,Q,c\) be positive integers with
\[
 (P,Q)=1,\quad P,Q\text{ odd},\quad
 5\mid Q,\quad5\nmid P,\quad P^2-5Q^2=4c^5.
\]
Then there are positive coprime odd integers \(t,s\), with \(5\nmid t\),
such that
\[
 16P=tG(t,s),\qquad16Q=5sF(t,s).
 \tag{Q4}\label{eq:quintic-half}
\]
\end{lemma}

\begin{proof}
Take
\(\alpha=(P+Q\sqrt5)/2=(P-Q)/2+Q\varphi\).
Its norm is \(c^5\), and its trace is \(P\).
The gcd of these two integers is \(1\), because a prime dividing both
would divide \(P\) and \(5Q^2\); also \(P\) is odd.
The product-and-sum Bezout argument again proves that
\(\alpha,\bar\alpha\) are coprime.  Its second coordinate is divisible
by \(5\) and its first is not.  Thus (Q1) gives \(\alpha=\gamma^5\).

Write \(\gamma=a+b\varphi=(t+s\sqrt5)/2\), where
\(s=b\), \(t=2a+b\).  The modulo-two calculation in the preceding
proof shows that \(b\) is odd, since the second coordinate of
\(\alpha\) is \(Q\).  Hence \(t,s\) are odd.
Multiplying \(\alpha=\gamma^5\) by \(32\) and applying (Q2) gives
(Q4).  Positivity of \(P,Q\) gives positivity of \(t,s\), and the
first coefficient modulo \(5\) gives \(5\nmid t\).
If an odd prime divided both \(t,s\), then \(\gamma\) and
\(\bar\gamma\) would both be divisible by that prime in \(R\);
the quotients remain in \(R\), because division of both odd
coordinates by an odd integer preserves their common parity.
Coprimality of the conjugates excludes this.  The prime \(2\) is
excluded because \(t,s\) are odd.
\end{proof}

The witnesses in these two lemmas can, if desired, be selected directly
by searching
\[
 1\le t,s\le16(P+Q+1).
 \tag{Q5}\label{eq:quintic-root-bound}
\]
The coefficient equations just proved place the witnesses inside this
box.  Thus no unbounded search for an algebraic fifth root occurs in
the actual descent.

\subsection{Two impossible families}

For \(h\ge0\), define an \emph{odd state} to be positive integers
\(t,s,w\) satisfying
\[
 \begin{gathered}
 (t,s)=1,\quad t,s\text{ odd},\quad5\mid s,\quad5\nmid t,\\
 5^h sF(t,s)=16w^5.
 \end{gathered}
 \tag{Q6}\label{eq:quintic-oddstate}
\]
Define an \emph{even state} by
\[
 \begin{gathered}
 (t,s)=1,\quad t\text{ odd},\quad s\text{ even},
 \quad5\mid s,\quad5\nmid t,\\
 2\cdot5^h sF(t,s)=w^5.
 \end{gathered}
 \tag{Q7}\label{eq:quintic-evenstate}
\]
The coefficient \(2\) in the even family stays fixed throughout the
descent; no additional parameter for its exponent is necessary.

The gcd calculations used below are
\[
 (s,F(t,s))=1,\qquad (5,F(t,s))=1.
 \tag{Q8}\label{eq:quintic-gcd}
\]
Indeed \(F(t,s)\equiv t^4\pmod s\) and
\(F(t,s)\equiv t^4\pmod5\).
For odd \(t,s\), \(16\mid F(t,s)\), since odd fourth powers are
\(1\pmod {16}\) and odd squares are \(1\) or \(9\pmod {16}\).
If \(t\) is odd and \(s\) even, \(F(t,s)\) is odd.

In an odd state, divide \(F\) by \(16\) and use (Q8).
The two coprime positive factors
\(5^h s\) and \(F(t,s)/16\) have product \(w^5\); therefore
\[
 5^h s=u^5,\qquad F(t,s)=16v^5
 \tag{Q9}\label{eq:quintic-oddsplit}
\]
for positive integers \(u,v\).
In an even state the corresponding conclusion is
\[
 2\cdot5^h s=u^5,\qquad F(t,s)=v^5.
 \tag{Q10}\label{eq:quintic-evensplit}
\]
These are ordinary integer prime-multiplicity arguments, with bounded
factorization or bounded searches for \(u,v\).

Now use the identity
\[
 (t^2+5s^2)^2-5(2s^2)^2=F(t,s).
 \tag{Q11}\label{eq:quintic-normidentity}
\]
For an odd state set \(P=(t^2+5s^2)/2\), \(Q=s^2\).
Both are odd, their gcd is \(1\), \(5\mid Q\), \(5\nmid P\), and
\(P^2-5Q^2=4v^5\).
Here \((P,s^2)=1\) follows from \((t,s)=1\), and the other claims
follow directly by reduction modulo \(4\) and \(5\).
Lemma~\ref{lem:quintic-half} produces positive coprime odd \(t',s'\)
with \(5\nmid t'\) and
\[
 16s^2=5s'F(t',s').
 \tag{Q12}\label{eq:quintic-oddnext}
\]
For an even state set \(P=t^2+5s^2\), \(Q=2s^2\).
Then \(P\) is odd, \(Q\) even, they are coprime, \(5\mid Q\),
\(5\nmid P\), and \(P^2-5Q^2=v^5\).
Lemma~\ref{lem:quintic-direct} gives positive coprime \(t',s'\),
\(t'\) odd, \(s'\) even, \(5\nmid t'\), and
\[
 2s^2=5s'F(t',s').
 \tag{Q13}\label{eq:quintic-evennext}
\]
In either case \(5\mid s'\).  The left side is divisible by \(25\);
cancelling one factor \(5\) on the right shows
\(5\mid s'F(t',s')\), while \(5\nmid F(t',s')\).

The new second coordinate is strictly smaller:
\[
 \boxed{s'<s.}
 \tag{Q14}\label{eq:quintic-decrease}
\]
In fact \(F(t',s')>5s'^4\), because \(t'>0\).
Writing either (Q12) or (Q13) as
\(Ks^2=5s'F(t',s')\), with \(K=16\) or \(2\), gives
\(25s'^5<Ks^2\).
If \(s'\ge s\), positivity gives
\(s^2\le s'^2\le s'^5\), contradicting \(K<25\).

Finally, square the first equation in (Q9) or (Q10).
For the odd state, (Q12) gives
\[
 5^{2h+1}s'F(t',s')
 =16(5^h s)^2
 =16(u^2)^5.
\]
For the even state, (Q13) gives
\[
 2\cdot5^{2h+1}s'F(t',s')
 =4\cdot5^{2h}s^2
 =(u^2)^5.
\]
Thus the successor is a state of the same type with
\[
 (h,t,s,w)\longmapsto(2h+1,t',s',u^2)
 \quad\hbox{and}\quad0<s'<s.
 \tag{Q15}\label{eq:quintic-transition}
\]

Neither state family has a member.  To see this in precisely the form
needed for iPRA, totalize the successor by returning a fixed default
on invalid input, and otherwise use the bounded searches described
above.  Starting from a purported state with second coordinate \(s_0\),
iterate this primitive-recursive function \(s_0\) times.
Quantifier-free induction on the iteration number proves that each
iterate remains a state and that
\(s_i+i\le s_0\).
At \(i=s_0\) this gives \(s_i\le0\), whereas every state has
\(s_i>0\).  This is a finite contradiction, not an invocation of
a well-ordering principle for arbitrary predicates.

\subsection{A fifth-power Fermat equation enters one of the families}

Suppose nonzero integers satisfy \(a^5+b^5=c^5\).
Cancelling their common gcd gives a pairwise coprime solution.

\paragraph{The modulo-25 preliminary.}
A fifth power modulo \(25\) is one of
\(0,1,7,-7,-1\); the nonzero four residues correspond exactly to
bases not divisible by \(5\).
No sum of three members of \(\{1,7,-7,-1\}\) is \(0\pmod {25}\).
For instance, their pairwise sums are
\(0,\pm2,\pm6,\pm8,\pm14\), none of which is
\(1,7,-7,-1\pmod {25}\).
Applying this to \(a^5+b^5+(-c)^5=0\) shows that one base is divisible
by \(5\).

Move that term to the right, changing signs if necessary, to obtain
\[
 x^5+y^5=5^5z^5,\qquad
 (x,y)=1,\qquad xyz\ne0,\qquad z>0.
 \tag{Q16}\label{eq:quintic-normalized}
\]
Neither \(x\) nor \(y\) is divisible by \(5\).
They cannot both be even.

\paragraph{Opposite parity.}
If \(x,y\) have opposite parity, put
\(p=x+y\), \(q=x-y\).  Then \(p,q\) are odd,
\((p,q)=1\), and
\[
 p(p^4+10p^2q^2+5q^4)=16\cdot5^5z^5.
\]
The quartic factor is positive, so \(p>0\).
Modulo \(5\) the left side is \(p^5\), hence \(p=5r\) with \(r>0\).
Coprimality implies \(5\nmid q\) and \((q,r)=1\);
in particular \(q\ne0\).
Writing \(q_0=|q|\), we get positive coprime odd \(q_0,r\), with
\(5\nmid q_0\), and
\[
 rG(q_0,r)=16\cdot5^3z^5.
 \tag{Q17}\label{eq:quintic-core}
\]

\paragraph{Both odd.}
If \(x,y\) are odd, put
\(p=(x+y)/2\), \(q=(x-y)/2\).
These are coprime and have opposite parity, and
\[
 2p(p^4+10p^2q^2+5q^4)=5^5z^5.
\]
Again \(p>0\) and \(p=5r\) with \(r>0\).
The right side is even, so \(z=2z_0\) with \(z_0>0\).
Substitution and cancellation give exactly
\[
 rG(|q|,r)=16\cdot5^3z_0^5.
\]
If \(r\) were odd, \(q\) would be even, and the left side would be
odd, a contradiction.  Thus \(r\) is even and \(|q|\) odd.
Again \((|q|,r)=1\), \(5\nmid q\), and \(q\ne0\).

We have therefore reached the common core
\[
 rG(q,r)=16\cdot5^3Z^5,\qquad
 q,r,Z>0,\quad (q,r)=1,\quad q\text{ odd},\quad5\nmid q,
 \tag{Q18}\label{eq:quintic-unified-core}
\]
where \(r\) is either odd or even.
We now spell out the last allocation of coprime factors.

The factors \(r,G(q,r)\) are coprime, since
\(G(q,r)\equiv q^4\pmod r\), and \(5\nmid G(q,r)\).
If \(r\) is odd, \(16\mid G(q,r)\); this follows from
odd fourth powers modulo \(16\) and
\(50\equiv2,\ 125\equiv13\pmod {16}\).
Thus
\[
 (25r)\frac{G(q,r)}{16}=(5Z)^5
\]
is a product of coprime positive integers, and so
\[
 25r=A^5,\qquad G(q,r)=16B^5.
 \tag{Q19}\label{eq:quintic-initial-odd-split}
\]
If \(r\) is even, \(G(q,r)\) is odd; hence
\[
 (50r)G(q,r)=(10Z)^5
\]
is a coprime factorization and
\[
 50r=A^5,\qquad G(q,r)=B^5.
 \tag{Q20}\label{eq:quintic-initial-even-split}
\]
Both cases imply \(5\mid r\): in the first,
\(5\mid A\) and \(r=125(A/5)^5\); in the second,
\(2r=125(A/5)^5\), which has the same consequence.

Use the second norm identity
\[
 G(q,r)=(q^2+25r^2)^2-5(10r^2)^2.
 \tag{Q21}\label{eq:quintic-initial-norm}
\]
For odd \(r\), take
\[
 P=(q^2+25r^2)/2,\qquad Q=5r^2.
\]
Then \(P,Q\) are odd, positive, and coprime,
\(5\mid Q\), \(5\nmid P\), and
\(P^2-5Q^2=4B^5\).
The coprimality is immediate from \((q,r)=1\), \(5\nmid q\),
and the oddness of \(P\).
Lemma~\ref{lem:quintic-half} gives positive coprime odd \(t,s\),
\(5\nmid t\), with
\[
 16r^2=sF(t,s).
\]
Since \(5\mid r\) and \(F(t,s)\equiv t^4\pmod5\), we have \(5\mid s\).
Squaring \(25r=A^5\) proves
\[
 5^4sF(t,s)=16(A^2)^5.
\]
This is an odd state with \(h=4\).

For even \(r\), take
\(P=q^2+25r^2\), \(Q=10r^2\).
These are coprime with \(P\) odd and \(Q\) even;
\(5\mid Q\), \(5\nmid P\), and
\(P^2-5Q^2=B^5\).
Lemma~\ref{lem:quintic-direct} gives positive coprime \(t,s\),
\(t\) odd, \(s\) even, \(5\nmid t\), and
\[
 2r^2=sF(t,s).
\]
Again \(5\mid s\).  Squaring \(50r=A^5\) yields
\[
 2\cdot5^4sF(t,s)=(A^2)^5,
\]
an even state with \(h=4\).

Both families were proved impossible.  This contradiction establishes:
\[
 \boxed{\forall a,b,c>0,\qquad a^5+b^5\ne c^5.}
 \tag{Q22}\label{eq:flt-five}
\]

\subsection{Primitive-recursive and Lean audit for the quintic proof}

All searches in the proof have explicit finite bounds.
Integer factorization can use trial division through the input.
The Euclidean algorithm in \(R\) has at most one more iteration than
the initial absolute norm.  Normalization of an element with coordinate
height \(H\) searches signed exponents between \(-4H\) and \(4H\).
Factorization in \(R\) searches divisor coordinates bounded by \(3M+6\),
where \(M\) is its absolute norm.  Algebraic-root extraction used by
the descent can instead search the box (Q5).
A fifth root of a positive integer \(D\) is sought only in \([1,D]\).
The main descent is iterated exactly \(s_0\) times.

For an especially simple envelope on the main states, if their four
coordinates \(h,t,s,w\) are at most \(B\ge1\), the next coordinates
are bounded by
\[
 T(B)=2^{100(B+1)}.
\]
Indeed \(h'=2h+1\), the two norm coordinates are bounded by
\(6(B+1)^2\), (Q5) bounds \(t',s'\) polynomially in \(B\), and
\(w'=u^2\le(2\cdot5^h s)^2\).
Iterating \(T\) exactly \(s_0\) times supplies one primitive-recursive
bound for the complete main-state trace.
For the internal Euclidean calculations, one may bound a coordinate
height after each division by
\(E(H)=32(H+1)^3\), and iterate \(E\) for the already bounded
number of divisions.  Efficiency is irrelevant; these bounds certify
primitive recursiveness.

A Lean completion should use integer-pair code for \(R\), not
\texttt{Real.sqrt}, and should isolate the following interfaces:
\[
\begin{array}{ll}
\texttt{goldenEuclideanStep}&
  \text{explicit quotient, remainder, and norm decrease},\\
\texttt{goldenUnitNormalForm}&
  u=\pm\varphi^j\text{ with a bounded signed exponent},\\
\texttt{goldenFifthExtraction}&
  \text{the coprime extraction implication (Q1)},\\
\texttt{fifthCoordinatesDirect}&
  \text{the direct bounded coordinate lemma},\\
\texttt{fifthCoordinatesHalf}&
  \text{the half-integral coordinate lemma},\\
\texttt{fifthOddStep},\ \texttt{fifthEvenStep}&
  \text{the transitions (Q15)},\\
\texttt{fifthInitialState}&
  \text{the finite normalization (Q16)--(Q21)}.
\end{array}
\]
The state predicates are decidable combinations of integer equalities,
inequalities, gcd, parity, and remainders.
Each interface should be proved as a quantifier-free correctness
statement about a specified primitive-recursive selector.
A final quantifier-free invariant on the bounded trace gives the iPRA
derivation.  Proving the same mathematical implication in Lean's
ambient logic is a separate, weaker deliverable than exhibiting that
object-theoretic iPRA derivation.

For orientation, the independently available source tree
\texttt{isomorphic-ai/fermat}, under
\texttt{Fermat/Five/}, separates the historical argument into
\texttt{Coordinates}, \texttt{PowerExtraction}, \texttt{Initial}, and
\texttt{Descent} modules.  It was inspected as a comparison source,
not compiled as part of this package.  Its ambient use of real embeddings,
choice, or least-counterexample formulations must not be counted as
an iPRA formalization.  The explicit bounded normalization and fixed-length
iteration above describe the replacements required here.

